\SourceLeaf{42}

\begin{center}
\textbf{NACHLASS.}
\end{center}

\begin{center}
\textbf{DE INTEGRATIONE FORMULAE DIFFERENTIALIS}\\
$(1+n\cos\varphi)^{\nu}. d\varphi.$
\end{center}

\noindent\hspace*{\parindent}\textls{Prooemium.} Maxima pars Calculi, qui Integralis vocatur, serierum evolutione et
tractatione absolvitur. Unica nimirum methodus, quam in illo praeter tentamina et
differentiatione in memoriam revocata invenies, hoc agit, ut data quantitas differentialis
in seriem expandatur terminorum, quorum singulus quilibet facile ad integrationem
perducitur. In serie finita hoc modo ad expressionem omnibus numeris absolutam
pervenitur; series infinita, licet incognitarum relationem quaesitam minus perfecte
exhibeat, modo ita procedat, ut adproximationi locus detur, id quod plurimis casibus apta
tractatione obtinebitur, ad vulgares usus sufficere et aequationi, si qua talis revera
inter variabiles propositas detur, aequipollere recte censetur. Est autem plerumque
propter terminorum seriem ingredientium prolixitatem et adparentem anomaliam arduum sane
negotium, legem istarum expressione generali circumscribere, circumscriptam
demonstratione munire; utrumque imminens erroris, irrepentis ex sola coniectura producto
calculo ad verisimilitatem elata, periculum et scientiae dignitas postulant. Exemplum
eiusmodi disquisitionis sistet dissertatiuncula proposita, quae in formulae
$d\varphi(1+n\cos\varphi)^{\nu}$, existentibus $n$ et $\nu$ numeris definitis pro lubitu
adsumtis, integratione per series absolvenda versatur. Formula haec, cuius eximius est in
Astronomia physica usus, adeo gravis Ill. \textsc{Eulero} visa fuit, ut illi integrum
Institutionum Calculi Integralis caput\HistoricalNote{*)}{*)}{Cap. VI T. I.} di%
